\documentclass{article}

\usepackage[utf8]{inputenc}
\usepackage{amsmath}
\usepackage{amssymb}
\usepackage{amsthm}
\usepackage{array}
\usepackage{geometry}
\usepackage{hyperref}
\usepackage{parskip}
\usepackage{setspace}

\geometry{
 a4paper,
 left=35mm,
 right=35mm,
 top=30mm,
}

\theoremstyle{plain}
\newtheorem{theorem}{Theorem}
\newtheorem{proposition}{Proposition}
\newtheorem{lemma}{Lemma}
\newtheorem{corollary}{Corollary}

\theoremstyle{definition}
\newtheorem{definition}{Definition}
\newtheorem{assumption}{Assumption}

\newcommand{\E}{\mathbb{E}}
\newcommand{\N}{\mathcal{N}}
\newcommand{\ind}{\mathbf{1}}
\newcommand{\doi}[1]{\href{https://doi.org/#1}{doi:#1}}

\title{Optimal enfranchisement}
\author{Ulysse Lojkine \and Claude Opus 4.7 \and Codex GPT-5.5}
\date{May 2026}

\onehalfspacing

\begin{document}

\maketitle

\begin{abstract}
A planner picks the franchise --- the subset of agents allowed to vote --- for a binary decision that affects different agents with different intensities.  This note revisits the optimal enfranchisement model of Fleurbaey and Lojkine~\cite{fleurbaey-lojkine} without their symmetry assumption on opinions.  Two natural conclusions from the symmetric model fail in full generality, even when the joint opinion distribution has full support: a planner's dictatorship can beat every majority franchise, and the unique optimal franchise need not be composed of the highest-stake agents.  We give two five-agent counterexamples.
\end{abstract}

\section{Introduction}

Many collective decisions are binary --- admit a country, lift a curfew, ban a substance --- and bear unequally on the population.  A normative theory that takes preference intensity seriously must decide who should vote.  We give a stylized model of this trade-off and ask two natural questions:
\begin{enumerate}
    \item Is voting always at least as good as letting the planner pick an outcome ex ante?
    \item Is the optimal franchise the one composed of the highest-stake agents?
\end{enumerate}
Fleurbaey and Lojkine~\cite{fleurbaey-lojkine} prove that, under ex ante permutation invariance of opinions, universal franchise beats every planner's dictatorship and every optimal franchise is high-stake.  They also analyze the common-iid case and the fair-coin prefix problem.  This note shows why their symmetry assumption is doing real work.  Under full support of the joint opinion distribution alone, the two positive conclusions need not survive: Section~\ref{sec:negative} gives a five-agent example where every odd majority franchise is beaten by the constant rule $Y \equiv 1$, and a second five-agent example where the unique optimal franchise excludes a strictly higher-stake agent in favour of a lower-stake one.

\paragraph{Related literature.}
This note is a companion to Fleurbaey and Lojkine~\cite{fleurbaey-lojkine}.  The closest classical results are the Condorcet jury theorems \cite{condorcet-sep}, which give independence and competence conditions under which majority correctness rises with the size of a homogeneous jury.  The expert-voting literature \cite{nitzan-paroush, berend} solves for optimal weighted majority rules when voter competences are observed, with weights $\log(p_i/(1-p_i))$.  Schmitz and Tr\"oger \cite{schmitz-troger} study ex-ante utilitarian welfare for binary alternatives with uncertain preference intensities; their main message agrees with the one here --- symmetry and neutrality make majority attractive, independence restores strong optimality, and correlated utilities can make majority suboptimal.  The Penrose--Banzhaf voting-power tradition \cite{penrose, banzhaf} stresses that a voter's marginal influence is governed by the probability of a tie among the others.  For schemes that aggregate intensity directly rather than choosing a franchise, see Casella and Gelman \cite{casella-gelman}.

\section{Model}\label{sec:model}

A planner faces a population $\N = \{1, \ldots, N\}$ and a binary decision $\{0,1\}$.  From her vantage, agent $i$'s opinion is a Bernoulli random variable $X_i$.  The planner picks a \emph{franchise} $C \subseteq \N$ and the outcome is the strict majority within $C$:
\[
    Y_C = \ind_{\{S_C > |C|/2\}},
    \qquad
    S_C = \sum_{i \in C} X_i .
\]
Strict majority is unbiased only when $|C|$ is odd, so we assume throughout that $N$ is odd and consider only nonempty franchises of odd size.  Let $p(i, C) = P[X_i = Y_C]$.

Each agent has utility $U(i, C)$ that depends on her own success.

\begin{assumption}[Selfishness]\label{ass:selfish}
Conditional on $\ind_{\{X_i = Y_C\}}$, agent $i$ is indifferent to the success of others and to anyone's inclusion in the franchise.
\end{assumption}

Under Assumption~\ref{ass:selfish}, $\E[U(i, C) \mid X_i = Y_C]$ and $\E[U(i, C) \mid X_i \neq Y_C]$ are constants in $C$.  Define
\[
    \Delta_i = \E[U(i, C) \mid X_i = Y_C] - \E[U(i, C) \mid X_i \neq Y_C] \geq 0,
\]
the agent's \emph{stake}.  Label agents so $\Delta_1 \geq \Delta_2 \geq \cdots \geq \Delta_N \geq 0$ and drop the constant terms from utility.  Aggregate expected welfare under franchise $C$ is
\begin{equation}\label{eq:welfare}
    U(C) = \sum_{i \in \N} \Delta_i\, p(i, C).
\end{equation}

\begin{assumption}[Full support]\label{ass:full-support}
For every $x \in \{0,1\}^N$, $P[(X_1, \ldots, X_N) = x] > 0$.
\end{assumption}

For comparison we consider a \emph{planner's dictatorship}: an outcome $Z_q \sim \mathrm{Bernoulli}(q)$ independent of $(X_1, \ldots, X_N)$, giving welfare $V(q) = \sum_{i \in \N} \Delta_i\, P[X_i = Z_q]$.  The objective is linear in $q$, so the best dictatorship is either $q = 0$ or $q = 1$.

\begin{definition}[High-stake franchise]\label{def:high-stake}
$C \subseteq \N$ is \emph{high-stake} if $C = \N$ or $\Delta_i \geq \Delta_j$ for every $i \in C$ and every $j \in \N \setminus C$.
\end{definition}

\section{Two negative results}\label{sec:negative}

\subsection{Dictatorship can beat every franchise}

\begin{proposition}\label{prop:dict}
There exist stakes $(\Delta_1, \ldots, \Delta_5)$ and a full-support distribution under which every odd franchise has strictly lower expected welfare than the dictatorship $q = 1$.
\end{proposition}

\begin{proof}
Take $N = 5$ and $(\Delta_1, \ldots, \Delta_5) = (5, 4, 3, 2, 1)$.  Define $P_0$ on $\{0,1\}^5$ by the following four atoms, with integer weights summing to $6$:
\[
\begin{array}{c|r}
x & 6\, P_0[X = x] \\
\hline
(0,1,1,1,0) & 1 \\
(1,0,0,1,1) & 3 \\
(1,0,1,0,1) & 1 \\
(1,1,0,0,0) & 1 .
\end{array}
\]
The dictatorship $q = 1$ has welfare $17/2$.  The welfares of all sixteen odd franchises, multiplied by $6$, are:
\[
\begin{array}{c|r@{\qquad}c|r}
C & 6\, U_{P_0}(C) & C & 6\, U_{P_0}(C) \\
\hline
\{1\}         & 48 & \{1,2,3\}     & 48 \\
\{2\}         & 45 & \{1,2,4\}     & 48 \\
\{3\}         & 45 & \{1,2,5\}     & 48 \\
\{4\}         & 45 & \{1,3,4\}     & 48 \\
\{5\}         & 45 & \{1,3,5\}     & 45 \\
\{1,4,5\}     & 45 & \{2,3,4\}     & 42 \\
\{2,3,5\}     & 45 & \{2,4,5\}     & 45 \\
\{3,4,5\}     & 48 & \{1,2,3,4,5\} & 48 .
\end{array}
\]
The best franchise welfare is $8$, so the dictatorship beats every franchise by at least $1/2$.

Mix $P_0$ with the uniform distribution $R$ on $\{0,1\}^5$ to enforce full support:
\[
    P = \tfrac{99}{100}\, P_0 + \tfrac{1}{100}\, R .
\]
Total stake is $15$, so replacing $P_0$ by $P$ can shrink any welfare gap by at most $15/100$.  Since
\[
    \tfrac{99}{100}\cdot\tfrac{1}{2} \;-\; \tfrac{15}{100}
    \;=\; \tfrac{69}{200} \;>\; 0,
\]
the dictatorship $q = 1$ still beats every franchise under the full-support distribution $P$.
\end{proof}

\subsection{Optimal franchises need not be high-stake}

\begin{proposition}\label{prop:high-stake}
There exist strictly decreasing stakes and a full-support distribution under which the unique optimal franchise is not high-stake.
\end{proposition}

\begin{proof}
Take $N = 5$ and $(\Delta_1, \ldots, \Delta_5) = (5, 4, 3, 2, 1)$.  Define $P_0$ on three states:
\[
\begin{array}{c|c}
x & P_0[X = x] \\
\hline
(1,0,0,0,1) & 1/5 \\
(1,0,0,1,1) & 3/5 \\
(1,1,0,0,0) & 1/5 .
\end{array}
\]
Let $C^\star = \{1, 2, 4\}$.  The welfares $5\, U_{P_0}(C)$ across the sixteen odd franchises are:
\[
\begin{array}{c|r@{\qquad}c|r}
C & 5\, U_{P_0}(C) & C & 5\, U_{P_0}(C) \\
\hline
\{1\}         & 39 & \{1,2,3\}     & 39 \\
\{2\}         & 39 & \{1,2,4\}     & 42 \\
\{3\}         & 36 & \{1,2,5\}     & 39 \\
\{4\}         & 39 & \{1,3,4\}     & 39 \\
\{5\}         & 36 & \{1,3,5\}     & 36 \\
\{1,4,5\}     & 36 & \{2,3,4\}     & 36 \\
\{2,3,5\}     & 36 & \{2,4,5\}     & 39 \\
\{3,4,5\}     & 39 & \{1,2,3,4,5\} & 39 .
\end{array}
\]
$C^\star$ is the unique maximizer, beating every other franchise by at least $3/5$.

Mixing with the uniform $R$ on $\{0,1\}^5$,
\[
    P = \tfrac{99}{100}\, P_0 + \tfrac{1}{100}\, R ,
\]
gives full support.  Total stake is $15$, so for any other franchise $D$,
\[
    U_P(C^\star) - U_P(D) \;\geq\; \tfrac{99}{100}\cdot\tfrac{3}{5} \;-\; \tfrac{1}{100}\cdot 15 \;=\; \tfrac{111}{250} \;>\; 0,
\]
and $C^\star$ remains the unique optimum.  Because the stakes are strictly decreasing, the only high-stake franchises of odd size are $\{1\}$, $\{1,2,3\}$, and $\{1,2,3,4,5\}$.  Since $C^\star$ excludes agent $3$ (stake $3$) while including agent $4$ (stake $2$), $C^\star$ is not high-stake.
\end{proof}

\paragraph{Intuition.}
In Proposition~\ref{prop:high-stake}, agent $3$'s opinion is constantly $0$ under $P_0$ and so conveys no information, while agent $4$'s opinion marks the high-probability state in which outcome $1$ best serves the high-stake agents.  Swapping voter $3$ for voter $4$ in the size-three franchise changes $Y$ in exactly the state $(1,0,0,1,1)$, flipping it from $0$ (pleasing stakes $4+3 = 7$) to $1$ (pleasing stakes $5+2+1 = 8$); the net gain $1 \times 3/5$ is the entire margin.  In Proposition~\ref{prop:dict}, the joint distribution is biased towards $X_1 = 1$, and the constant dictator who sides with the highest-stake agent outperforms any majority of a fluctuating subset.

\section{Relation to existing positive results}\label{sec:positive}

The two counterexamples above are outside the symmetric-opinions domain analyzed by Fleurbaey and Lojkine~\cite{fleurbaey-lojkine}.  In that paper, full support plus ex ante permutation invariance of $(X_i)_{i \in \N}$ is enough to prove that universal franchise has strictly larger expected aggregate welfare than any planner's dictatorship, and that every optimal franchise is high-stake.  Their notation writes the success probabilities as $p_+(c)$ for insiders and $p_-(c)$ for outsiders in a franchise of size $c$; their key lemma is $p_+(c) > p_-(c)$, which immediately yields high-stake optimality by exchanging a lower-stake insider for a higher-stake outsider.

The iid cases that one might use as positive benchmarks here are therefore already covered there.  When opinions are common iid Bernoulli, Fleurbaey and Lojkine derive the fixed-size comparison criterion and the associated one-dimensional search over high-stake prefixes; in the balanced case $p=1/2$, this gives the familiar central-binomial term governing whether a larger prefix is worthwhile.  This note does not reprove those results.  Its purpose is instead to show that, once permutation invariance is dropped, full support alone is too weak: the dictatorship comparison and the high-stake conclusion can both fail.

\begin{thebibliography}{9}

\bibitem{fleurbaey-lojkine}
Marc Fleurbaey and Ulysse Lojkine.
\newblock Optimal enfranchisement.
\newblock \emph{Social Choice and Welfare}, 2025.
\newblock \doi{10.1007/s00355-025-01607-9}.
\newblock \url{https://shs.hal.science/halshs-05158588}.

\bibitem{condorcet-sep}
Franz Dietrich and Kai Spiekermann.
\newblock Jury Theorems.
\newblock \emph{Stanford Encyclopedia of Philosophy}, 2021.
\newblock \url{https://plato.stanford.edu/entries/jury-theorems/}.

\bibitem{nitzan-paroush}
Shmuel Nitzan and Jacob Paroush.
\newblock Optimal decision rules in uncertain dichotomous choice situations.
\newblock \emph{International Economic Review}, 23(2):289--297, 1982.

\bibitem{berend}
Daniel Berend, Luba Sapir, and Shmuel Sapir.
\newblock Ranking of weighted majority rules.
\newblock \emph{Journal of Applied Probability}, 45(4):994--1006, 2008.

\bibitem{schmitz-troger}
Patrick W.\ Schmitz and Thomas Tr\"oger.
\newblock The (sub-)optimality of the majority rule.
\newblock \emph{Games and Economic Behavior}, 74(2):651--665, 2012.

\bibitem{penrose}
L.\ S.\ Penrose.
\newblock The elementary statistics of majority voting.
\newblock \emph{Journal of the Royal Statistical Society}, 109:53--57, 1946.

\bibitem{banzhaf}
John F.\ Banzhaf III.
\newblock Weighted voting doesn't work: a mathematical analysis.
\newblock \emph{Rutgers Law Review}, 19:317--343, 1965.

\bibitem{casella-gelman}
Alessandra Casella and Andrew Gelman.
\newblock A simple scheme to improve the efficiency of referenda.
\newblock \emph{Journal of Public Economics}, 92(10--11):2240--2261, 2008.

\end{thebibliography}

\end{document}
